\documentclass[aspectratio=169]{beamer}

\usetheme{Madrid}

\title{La rédaction scientifique avec \LaTeX}
\subtitle{Opportunité pour les chercheurs/auteurs et les revues scientifiques}
\author{Dr Akpoué Josué}
\date{13 mai 2026}


% ---------------------------------------------------
% CHARGEMENT DES PAQUETS
% ---------------------------------------------------

\usepackage[utf8]{inputenc}
\usepackage{linguex,tikz-qtree,tipa}
\usepackage{verbatim}



\begin{document}
	
	%------------------------------------------------
	\begin{frame}
		\titlepage
	\end{frame}
	%------------------------------------------------
	
	\begin{frame}{Introduction}
		\begin{itemize}
			\item La publication scientifique exige aujourd’hui des documents de plus en plus rigoureux sur les plans :
			\begin{itemize}
				\item scientifique ;
				\item typographique ;
				\item éditorial.
			\end{itemize}
			
			\vspace{0.3cm}
			
			\item Les chercheurs en linguistique produisent souvent des documents complexes :
			\begin{itemize}
				\item transcriptions phonétiques ;
				\item exemples glossés ;
				\item tableaux ;
				\item arbres syntaxiques ;
				\item bibliographies importantes.
			\end{itemize}
			
			\vspace{0.3cm}
			
			\item Les traitements de texte classiques montrent rapidement leurs limites :
			\begin{itemize}
				\item incohérences de mise en page ;
				\item gestion difficile des références ;
				\item instabilité des documents longs ;
				\item travail éditorial lourd pour les revues.
			\end{itemize}
			
			\vspace{0.3cm}
			
			\item \LaTeX{} propose une approche différente :
			\begin{itemize}
				\item séparation du contenu et de la mise en forme ;
				\item automatisation de nombreuses tâches ;
				\item production de documents scientifiques de qualité professionnelle.
			\end{itemize}
		\end{itemize}
	\end{frame}
	
	%------------------------------------------------
	
	\begin{frame}{Objectifs de la présentation}
		\begin{itemize}
			\item Comprendre ce qu’est \LaTeX{} et son origine ;
			
			\vspace{0.2cm}
			
			\item Découvrir la philosophie de rédaction scientifique basée sur la structure logique des documents ;
			
			\vspace{0.2cm}
			
			\item Présenter quelques usages de \LaTeX{} dans les travaux linguistiques ;
			
			\vspace{0.2cm}
			
			\item Montrer les avantages de \LaTeX{} :
			\begin{itemize}
				\item pour les auteurs ;
				\item pour les directeurs de revue ;
				\item pour les équipes éditoriales.
			\end{itemize}
			
			\vspace{0.2cm}
			
			\item Introduire les bases nécessaires pour :
			\begin{itemize}
				\item rédiger un manuscrit ;
				\item créer des modèles ;
				\item automatiser la mise en page d’une revue scientifique.
			\end{itemize}
		\end{itemize}
	\end{frame}
	
	%------------------------------------------------
	
	%------------------------------------------------
	\section{Qu’est-ce que \LaTeX{} ?}
	%------------------------------------------------
	
	\begin{frame}{Qu’est-ce que \LaTeX{} ?}
		\begin{itemize}
			\item \LaTeX{} est un système de composition de documents scientifiques.
			
			\vspace{0.3cm}
			
			\item Il permet de produire des documents de haute qualité typographique :
			\begin{itemize}
				\item articles ;
				\item livres ;
				\item mémoires ;
				\item présentations ;
				\item posters scientifiques.
			\end{itemize}
			
			\vspace{0.3cm}
			
			\item Contrairement aux traitements de texte classiques, \LaTeX{} fonctionne à partir :
			\begin{itemize}
				\item d’un fichier texte ;
				\item d’instructions de structuration ;
				\item d’une compilation vers un document final (PDF).
			\end{itemize}
			
			\vspace{0.3cm}
			
			\item L’utilisateur décrit la structure logique du document :
			\begin{itemize}
				\item titre ;
				\item section ;
				\item citation ;
				\item bibliographie ;
				\item exemple linguistique ;
				\item figure ;
				\item tableau.
			\end{itemize}
		\end{itemize}
	\end{frame}
	
	%------------------------------------------------
	
	\begin{frame}{De TeX à \LaTeX{} : bref historique}
		\begin{itemize}
			\item À la fin des années 1970, Donald Knuth développe TeX.
			
			\vspace{0.3cm}
			
			\item Motivation initiale :
			\begin{itemize}
				\item produire des ouvrages scientifiques de très haute qualité ;
				\item améliorer la composition mathématique ;
				\item garantir une typographie stable et professionnelle.
			\end{itemize}
			
			\vspace{0.3cm}
			
			\item TeX est puissant, mais relativement complexe à utiliser.
			
			\vspace{0.3cm}
			
			\item Leslie Lamport crée ensuite \LaTeX{} :
			\begin{itemize}
				\item simplification de TeX ;
				\item introduction de commandes plus accessibles ;
				\item automatisation de nombreuses tâches de mise en page.
			\end{itemize}
			
			\vspace{0.3cm}
			
			\item Aujourd’hui, \LaTeX{} est largement utilisé dans le monde académique et scientifique.
		\end{itemize}
	\end{frame}
	
	%------------------------------------------------
	
	\begin{frame}{WYSIWYG vs WYSIWYM}
		\begin{columns}[T]
			
			\column{0.48\textwidth}
			
			\textbf{WYSIWYG}
			
			\vspace{0.2cm}
			
			\small
			\emph{What You See Is What You Get}
			
			\vspace{0.3cm}
			
			\begin{itemize}
				\item Mise en forme directe ;
				\item Approche visuelle ;
				\item Exemple :
				\begin{itemize}
					\item Microsoft Word ;
					\item LibreOffice Writer.
				\end{itemize}
				
				\vspace{0.2cm}
				
				\item L’utilisateur contrôle manuellement :
				\begin{itemize}
					\item polices ;
					\item tailles ;
					\item espacements ;
					\item alignements.
				\end{itemize}
			\end{itemize}
			
			\column{0.48\textwidth}
			
			\textbf{WYSIWYM}
			
			\vspace{0.2cm}
			
			\small
			\emph{What You See Is What You Mean}
			
			\vspace{0.3cm}
			
			\begin{itemize}
				\item Approche structurelle ;
				\item Séparation contenu / présentation ;
				\item Exemple :
				\begin{itemize}
					\item \LaTeX{}
				\end{itemize}
				
				\vspace{0.2cm}
				
				\item L’utilisateur décrit :
				\begin{itemize}
					\item la structure ;
					\item le rôle logique ;
					\item la sémantique du contenu.
				\end{itemize}
			\end{itemize}
			
		\end{columns}
		
		\vspace{0.4cm}
		
		\begin{block}{Conséquence}
			Cette approche favorise la cohérence, l’automatisation et la stabilité des documents scientifiques.
		\end{block}
		
	\end{frame}
	
	%------------------------------------------------
	
	%------------------------------------------------
	\section{Applications de \LaTeX{}}
	%------------------------------------------------
	
	\begin{frame}{Les principales applications de \LaTeX{}}
		\begin{itemize}
			\item Rédaction d’articles scientifiques ;
			
			\vspace{0.2cm}
			
			\item Production de livres, mémoires et thèses ;
			
			\vspace{0.2cm}
			
			\item Création de présentations académiques avec Beamer ;
			
			\vspace{0.2cm}
			
			\item Réalisation de posters scientifiques ;
			
			\vspace{0.2cm}
			
			\item Production de documents techniques et professionnels ;
			
			\vspace{0.2cm}
			
			\item Automatisation de la mise en page éditoriale ;
			
			\vspace{0.2cm}
			
			\item Gestion avancée des références bibliographiques ;
			
			\vspace{0.2cm}
			
			\item Création de documents scientifiques multilingues.
		\end{itemize}
	\end{frame}
	
	%------------------------------------------------
	
	\begin{frame}{Applications en linguistique}
		\begin{itemize}
			\item \textbf{Transcriptions phonétiques}
			\begin{itemize}
				\item alphabet phonétique international (API/IPA) ;
				\item caractères Unicode ;
				\item diacritiques spécialisés.
			\end{itemize}
			
			\vspace{0.3cm}
			
			\item \textbf{Exemples linguistiques}
			\begin{itemize}
				\item numérotation automatique ;
				\item alignement cohérent ;
				\item références croisées.
			\end{itemize}
			
			\vspace{0.3cm}
			
			\item \textbf{Gloses interlinéaires}
			\begin{itemize}
				\item alignement automatique des morphèmes ;
				\item présentation normalisée.
			\end{itemize}
		\end{itemize}
	\end{frame}
	
	%------------------------------------------------
	
	\begin{frame}{Autres usages utiles en linguistique}
		\begin{itemize}
			\item \textbf{Arbres syntaxiques}
			\begin{itemize}
				\item représentation graphique des structures syntaxiques ;
				\item packages spécialisés.
			\end{itemize}
			
			\vspace{0.3cm}
			
			\item \textbf{Tableaux linguistiques}
			\begin{itemize}
				\item paradigmes morphologiques ;
				\item tableaux comparatifs ;
				\item données lexicales.
			\end{itemize}
			
			\vspace{0.3cm}
			
			\item \textbf{Bibliographies scientifiques}
			\begin{itemize}
				\item gestion automatique des références ;
				\item styles bibliographiques ;
				\item cohérence des citations.
			\end{itemize}
			
			\vspace{0.3cm}
			
			\item \textbf{Documents multilingues}
			\begin{itemize}
				\item coexistence de plusieurs langues ;
				\item prise en charge de différents systèmes d’écriture.
			\end{itemize}
		\end{itemize}
	\end{frame}
	
	%------------------------------------------------
	
	%------------------------------------------------
	\section{Pourquoi utiliser \LaTeX{} en linguistique ?}
	%------------------------------------------------
	
	\begin{frame}{Avantages pour les auteurs}
		\begin{itemize}
			\item \textbf{Qualité typographique professionnelle}
			\begin{itemize}
				\item gestion des espacements ;
				\item justification ;
				\item césures ;
				\item homogénéité du document.
			\end{itemize}
			
			\vspace{0.3cm}
			
			\item \textbf{Gestion automatique des références}
			\begin{itemize}
				\item citations ;
				\item bibliographie ;
				\item renvois internes ;
				\item numérotation automatique.
			\end{itemize}
			
			\vspace{0.3cm}
			
			\item \textbf{Adapté aux documents complexes}
			\begin{itemize}
				\item exemples glossés ;
				\item IPA ;
				\item arbres syntaxiques ;
				\item tableaux ;
				\item documents longs.
			\end{itemize}
			
			\vspace{0.3cm}
			
			\item \textbf{Gain de temps à long terme}
			\begin{itemize}
				\item automatisation ;
				\item réutilisation de modèles ;
				\item stabilité des documents.
			\end{itemize}
		\end{itemize}
	\end{frame}
	
	%------------------------------------------------
	
	\begin{frame}{Avantages pour les revues scientifiques}
		\begin{itemize}
			\item \textbf{Uniformisation des articles}
			\begin{itemize}
				\item même structure ;
				\item même typographie ;
				\item même style bibliographique.
			\end{itemize}
			
			\vspace{0.3cm}
			
			\item \textbf{Réduction du travail éditorial}
			\begin{itemize}
				\item moins de corrections manuelles ;
				\item automatisation de la mise en page ;
				\item cohérence des publications.
			\end{itemize}
			
			\vspace{0.3cm}
			
			\item \textbf{Création de modèles de revue}
			\begin{itemize}
				\item classes de documents ;
				\item fichiers de style ;
				\item modèles auteurs.
			\end{itemize}
			
			\vspace{0.3cm}
			
			\item \textbf{Professionnalisation de la publication}
			\begin{itemize}
				\item qualité comparable aux standards internationaux ;
				\item meilleure stabilité des archives numériques ;
				\item pérennité des documents.
			\end{itemize}
		\end{itemize}
	\end{frame}
	
	%------------------------------------------------
	
	\begin{frame}{Pourquoi \LaTeX{} est particulièrement pertinent en linguistique}
		\begin{block}{Une discipline à forte complexité typographique}
			Les travaux linguistiques combinent souvent :
			\begin{itemize}
				\item plusieurs langues ;
				\item caractères spéciaux ;
				\item transcriptions phonétiques ;
				\item gloses interlinéaires ;
				\item structures syntaxiques ;
				\item tableaux complexes.
			\end{itemize}
		\end{block}
		
		\vspace{0.4cm}
		
		\begin{alertblock}{Conséquence}
			\LaTeX{} permet de gérer cette complexité de manière plus stable, plus cohérente et plus professionnelle que les solutions de mise en page classiques.
		\end{alertblock}
		
	\end{frame}
	
	%------------------------------------------------
	
	
	%------------------------------------------------
	\section{Conclusion}
	%------------------------------------------------
	
	\begin{frame}{Conclusion}
		\begin{itemize}
			\item \LaTeX{} est bien plus qu’un simple outil de mise en page :
			\begin{itemize}
				\item c’est un système de composition scientifique ;
				\item une approche structurée de la rédaction académique ;
				\item un outil d’automatisation éditoriale.
			\end{itemize}
			
			\vspace{0.4cm}
			
			\item Pour les chercheurs et auteurs, il offre :
			\begin{itemize}
				\item une meilleure qualité typographique ;
				\item une gestion efficace des documents complexes ;
				\item une plus grande stabilité des manuscrits.
			\end{itemize}
			
			\vspace{0.4cm}
			
			\item Pour les revues scientifiques, il permet :
			\begin{itemize}
				\item l’uniformisation des publications ;
				\item la réduction du travail de mise en page ;
				\item la création de chaînes éditoriales plus professionnelles.
			\end{itemize}
			
			\vspace{0.4cm}
			
			\item Dans le domaine de la linguistique, où les exigences typographiques sont particulièrement élevées, \LaTeX{} constitue une solution puissante et durable.
		\end{itemize}
	\end{frame}
	
	%------------------------------------------------
	
	\begin{frame}{Perspectives}
		\begin{itemize}
			\item Développement de modèles adaptés aux revues locales ;
			
			\vspace{0.3cm}
			
			\item Création de fichiers de style pour automatiser la publication ;
			
			\vspace{0.3cm}
			
			\item Formation des chercheurs et étudiants à la rédaction scientifique structurée ;
			
			\vspace{0.3cm}
			
			\item Intégration progressive de \LaTeX{} dans les pratiques éditoriales des sciences du langage.
		\end{itemize}
		
		\vspace{0.6cm}
		
		\begin{center}
			\Large
			Merci pour votre attention
		\end{center}
		
	\end{frame}
	
	\begin{frame}
		\begin{center}
			\Large
			Merci pour votre attention
		\end{center}
	\end{frame}
	
	%------------------------------------------------
	
	\appendix
	
	%------------------------------------------------
	\section{Travaux dirigés}
	%------------------------------------------------
	
	\begin{frame}[fragile]{TD1 — Document minimal}
		\textbf{Objectif :} créer et compiler un premier document \LaTeX{}.
		
		\vspace{0.4cm}
		
		\textbf{Exemple :}
		
		\small
		
		\begin{verbatim}
			\documentclass{article}
 
			\begin{document}
				
				Bonjour le monde !
				
			\end{document}
		\end{verbatim}
		
		\normalsize
		
		\vspace{0.3cm}
		
		\textbf{Points abordés :}
		\begin{itemize}
			\item structure minimale ;
			\item compilation PDF ;
			\item rôle du préambule.
		\end{itemize}
		
	\end{frame}
	
	%------------------------------------------------
	
	\begin{frame}[fragile]{TD2 — Métadonnées et page de titre}
		\textbf{Objectif :} ajouter les informations du document.
		
		\vspace{0.4cm}
		
		\small
		
		\begin{verbatim}
			\title{Titre du document}
			\author{Nom de l'auteur}
			\date{\today}
			
			\begin{document}
				
				\maketitle
				
				\begin{abstract}
					Résumé du document.
				\end{abstract}
				
			\end{document}
		\end{verbatim}
		
		\normalsize
		
		\vspace{0.3cm}
		
		\textbf{Points abordés :}
		\begin{itemize}
			\item titre ;
			\item auteur ;
			\item date ;
			\item résumé.
		\end{itemize}
		
	\end{frame}
	
	%------------------------------------------------
	
	\begin{frame}[fragile]{TD3 — Structuration du document}
		\textbf{Objectif :} organiser un texte scientifique.
		
		\vspace{0.4cm}
		
		\small
		
		\begin{verbatim}
			\section{Introduction}
			
			\subsection{Contexte}
			
			\subsection{Objectifs}
			
			\section{Méthodologie}
		\end{verbatim}
		
		\normalsize
		
		\vspace{0.3cm}
		
		\textbf{Points abordés :}
		\begin{itemize}
			\item sections ;
			\item sous-sections ;
			\item hiérarchie logique du document.
		\end{itemize}
		
	\end{frame}
	
	%------------------------------------------------
	
	\begin{frame}[fragile]{TD4 — Mise en forme du texte}
		\textbf{Objectif :} appliquer des styles simples.
		
		\vspace{0.4cm}
		
		\small
		
		\begin{verbatim}
			\textit{Texte en italique}
			
			\textbf{Texte en gras}
			
			\emph{Texte emphatique}
			
			\textsc{Petites capitales}
		\end{verbatim}
		
		\normalsize
		
		\vspace{0.3cm}
		
		\textbf{Points abordés :}
		\begin{itemize}
			\item italiques ;
			\item gras ;
			\item emphase ;
			\item petites capitales.
		\end{itemize}
		
	\end{frame}
	
	%------------------------------------------------
	
	\begin{frame}[fragile]{TD5 — Références croisées et bibliographie}
		\textbf{Objectif :} automatiser les renvois et les références.
		
		\vspace{0.4cm}
		
		\small
		
		\begin{verbatim}
			\section{Introduction}
			\label{sec:intro}
			
			Voir la section \ref{sec:intro}.
			
			Selon \cite{dupont2020} ...
			
			\bibliographystyle{plain}
			\bibliography{references}
		\end{verbatim}
		
		\normalsize
		
		\vspace{-0.3cm}
		
		\textbf{Points abordés :}
		\begin{itemize}
			\item \verb|\label| ;
			\item \verb|\ref| ;
			\item \verb|\cite| ;
			\item bibliographie automatique.
		\end{itemize}
		
	\end{frame}
	
	%------------------------------------------------
	
	\begin{frame}[fragile]{TD6 — Transcription phonétique}
		\textbf{Objectif :} écrire des transcriptions phonétiques.
		
		\vspace{0.3cm}
		
		{Avec le paquet \verb|tipa| ou \verb|xsipa| :}
		
		\vspace{0.3cm}
		
		\small
		
		\begin{verbatim}
			\usepackage{tipa}
			
			[b\textepsilon\ng]
			
			\usepackage{xsipa}
			
			[b\xsE]
		\end{verbatim}
		
		
		\normalsize
		
		\vspace{0.3cm}
		
		\textbf{Points abordés :}
		\begin{itemize}
			\item Unicode ;
			\item caractères IPA ;
			\item polices linguistiques.
		\end{itemize}
		
	\end{frame}
	
	%------------------------------------------------
	
	\begin{frame}[fragile]{TD7 — Exemples linguistiques avec linguex}
		\textbf{Objectif :} produire des exemples linguistiques numérotés.
		
		\vspace{0.4cm}
		
		\small
		
		\begin{verbatim}
			\usepackage{linguex}
			
			\ex.
			Ceci est un exemple linguistique.
			
			\ex.
			\a. Première phrase
			\b. Deuxième phrase
		\end{verbatim}
		
		\normalsize
		
		\vspace{0.3cm}
		
		\textbf{Points abordés :}
		\begin{itemize}
			\item numérotation automatique ;
			\item sous-exemples ;
			\item alignement.
		\end{itemize}
		
	\end{frame}
	
	%------------------------------------------------
	
	\begin{frame}[fragile]{TD8 — Exemples glossés avec linguex}
		\textbf{Objectif :} créer des gloses interlinéaires.
		
		\vspace{0.4cm}
		
		\small
		
		\begin{verbatim}
			\usepackage{linguex}
			
			\exg.
			Kofi bE-RE ba\\
			Kofi venir-PROG ici\\
			`Kofi vient ici.'
		\end{verbatim}
		
		\normalsize
		
		\vspace{0.3cm}
		
		\textbf{Points abordés :}
		\begin{itemize}
			\item ligne source ;
			\item ligne de gloses ;
			\item traduction.
		\end{itemize}
		
	\end{frame}
	
	%------------------------------------------------
	
	\begin{frame}[fragile]{TD9 — Insertion d’images}
		\textbf{Objectif :} intégrer des figures dans un document.
		
		\vspace{0.4cm}
		
		\small
		
		\begin{verbatim}
			\usepackage{graphicx}
			
			\begin{figure}
				\centering
				\includegraphics[width=.6\textwidth]{image}
				\caption{Une image}
			\end{figure}
		\end{verbatim}
		
		\normalsize
		
		\vspace{0.3cm}
		
		\textbf{Points abordés :}
		\begin{itemize}
			\item insertion d’images ;
			\item taille ;
			\item légendes ;
			\item centrage.
		\end{itemize}
		
	\end{frame}
	
	%------------------------------------------------
	
	\begin{frame}[fragile]{TD10 — Arbres syntaxiques avec tikz-qtree}
		\textbf{Objectif :} produire des arbres syntaxiques.
		
		\vspace{0.4cm}
		
		\small
		
		\begin{verbatim}
			\usepackage{tikz-qtree}
			
			\begin{tikzpicture}
				\Tree [.TP
				[.DP Kofi ]
				[.T'
				[.T ]
				[\qroof{mange du riz}.VP ]]]
			\end{tikzpicture}
		\end{verbatim}
		
		\normalsize
		
		\vspace{0.3cm}
		
		\textbf{Points abordés :}
		\begin{itemize}
			\item structure hiérarchique ;
			\item nœuds syntaxiques ;
			\item représentation graphique.
		\end{itemize}
		
	\end{frame}
	
	%------------------------------------------------
	
	\begin{frame}[fragile]{TD11 — Fontspec et polices OpenType}
		\textbf{Objectif :} utiliser des polices externes avec XeLaTeX/LuaLaTeX.
		
		\vspace{0.4cm}
		
		\small
		
		\begin{verbatim}
			\usepackage{fontspec}
			
			\setmainfont{Times New Roman}
			
			\setsansfont{Arial}
			
			\setmonofont{JetBrains Mono}
		\end{verbatim}
		
		\normalsize
		
		\vspace{0.3cm}
		
		\textbf{Points abordés :}
		\begin{itemize}
			\item polices système ;
			\item OpenType ;
			\item TrueType ;
			\item compatibilité Unicode.
		\end{itemize}
		
	\end{frame}
	
	%------------------------------------------------
	
\end{document}